\section{はじめに}\label{sec:intro}

\subsection{動機}\label{sec:intro:motivation}

自律研究システムには、問いから実験で裏づけられた答えまでの距離を
縮める可能性がある。一方で、何を試すか、どの結果を信頼するか、何を
証拠として残すか、失敗からどう立て直すかという判断を避けることは
できない。提案役、評価役、敵対検査役、執筆役、裁定役は一つの研究制度
を形作り、その振る舞いの多くは指示文と採点規則で決まる。

制度を固定すれば理解しやすいが、初期設計の弱点も固定される。反対に、
制度を自由に更新できるようにすると、採点規則の変更者が自らに有利な
根拠を作り、査読役が自分の後継を選び、広い権限を得た指示文が不都合な
履歴を隠すおそれがある。したがって、自己改善は単なる最適化ではなく、
権限と説明責任の問題でもある。

\subsection{公開条件としての検証済み主張}\label{sec:intro:verified}

Constitutional ARI-RQGM（自律研究基盤 ARI の立憲型 RQGM）は、モデル
出力が制度変更を\emph{提案}する
ことはできても、直接\emph{実行}することはできない、という原則から
出発する。決定論的なカーネルが、記録形式、ハッシュ値、権限、役割
分離、遷移の適法性、監査履歴の連続性、隔離再生成の入力範囲を検査
する。制度台帳を書き換えられるのは一つの遷移処理だけである。制度の
状態は一つの期の間は固定され、境界での更新は全体が確定するか、再開時に
存在しなかったものとして扱われる。重大なカーネル違反に対する T16 は
現期を強制終了し、新しい指紋値の期を開始する。

固定層は、論文の構造化された数値主張と実行証拠を照合する最終判定も
保持する。統治された執筆役と査読役は原稿を比較し改稿できるが、宣言
済みの式と被演算値から数値を再計算する検証器は変更できない。一般の
自然言語主張の真偽や、実行記録自体の真実性は検証対象外である。

\subsection{貢献}\label{sec:intro:contrib}

本稿の貢献は次の五点である。

\begin{enumerate}
  \item 不変の憲法層、更新可能な制度層、後継を提案できても任命は
        できないメタ層からなる、三層の権限構造を定める。
  \item 指示文、構成要素、効用規則を一つの期の間固定し、
        境界更新と重大違反時の緊急境界を構成要素遷移表として示す。
  \item 成果物への攻撃から、弁護、裁定、検証済み攻撃までを分離し、
        作成時の期を記録した研究生成役と論文役へ責任を結ぶ。来歴が
        曖昧な場合は、対象を推測しない。
  \item 選択的な影響除去、隔離再生成、過去事例による再試験、来歴を
        考慮した修復により、退役を安全に反映する。効用規則を更新
        した場合は、保存済みの評価軸から再採点する。
  \item 同じ憲法を論文保管庫へ適用し、公開試験と固定故障例を報告し、
        未実施の対照実験を明示する。
\end{enumerate}

\subsection{本稿の範囲}\label{sec:intro:scope}

本稿の主題は制度の共進化であり、ARI の全実行機能を列挙することでは
ない。実験実行系、分野別の道具、成果物の組立て処理は、型の定まった成果物
を生成・消費する外部環境として扱う。内部の実行順序や探索方式は対象外
である。記録、来歴、保存点の契約を満たす研究実行系であれば、同じ
憲法上の主張が成立することを目標とする。

また、構成から保証できる性質と、実験で確かめるべき主張を区別する。
台帳の単一書き手、適法な遷移、局所的なハッシュ連鎖の検査、違法な状態
変更の拒否は構造上の性質である。外部固定のない末尾巻戻し、責任帰属の
意味的正しさ、新規性、研究品質、評価の校正、費用対効果には別の証拠が
要る。現状の結果はそれらを支持しない。

\subsection{構成}\label{sec:intro:organisation}

\Cref{sec:overview} でシステム境界と脅威を定める。
\Cref{sec:rqgm} で立憲型の共進化、期の境界、論文統治、監査記録を
説明する。\Cref{chapter:eval} で確認済みの挙動と未検証の主張を分け、
\Cref{sec:related} で自己改善・自律研究・評価研究との関係を整理する。
\Cref{sec:discussion} では固定憲法の代償と今後の課題を論じる。付録には
構成要素状態遷移表、権限行列、主要手続きと、RQGM 固有の指示文の目的を
まとめる。
